% Corrected French for printed p.208.
% Three print-level repairs and one page-seam normalization are recorded.
\clearpage
\SourcePage{208}

\LayerTarget{tohoku:proof:lemma-5.5.2:p1}
\noindent ctérisation des objets injectifs dans $\mathbf K({\mathbf C}_1)$ indiquée dans 2.4, $\mathrm C$ est
acyclique en dimensions $>0$ et se ``décompose'', de plus les $\mathrm C^i$ et $Z^i(\mathrm C)$
sont injectifs, d'où il résulte d'abord que $\overline F(\mathrm C)$ est une résolution
de $F(Z^0(\mathrm C))$ d'où ${\mathbf R}^nG(\overline F(\mathrm C))=R^nG(F(Z^0(\mathrm C)))$, et comme $Z^0(\mathrm C)$
est injectif le deuxième membre est nul d'après l'hypothèse sur $F$. Ainsi, on
peut appliquer le théorème 2.4.1 au foncteur composé
$(H^0\overline G)\overline F$, ce qui donne le foncteur spectral de l'énoncé,
compte tenu que $R^q\overline F=\overline{R^qF}$ (relation qui résulte
immédiatement des définitions).

\LayerTarget{tohoku:corollary:lemma-5.5.2}
\textsc{Corollaire.} \textit{Sous les conditions précédentes, soit $\mathrm C$
une résolution d'un objet $A\in {\mathbf C}_1$. Alors il existe une suite spectrale
cohomologique, aboutissant à l'objet gradué $(R^n(GF)(A))$, et dont le terme
initial est donné dans \textup{(5.5.2)}.}

\LayerTarget{tohoku:proof:corollary-lemma-5.5.2}
En effet, on a ici ${\mathbf R}^n(GF)(\mathrm C)=R^n(GF)(A)$.

\LayerTarget{tohoku:definition:g-cover}
Revenons à l'espace $X$ avec le groupe d'opérateurs $G$. Soit
${\mathbf U}=(U_i)_{i\in I}$ un recouvrement de $X$, $G$ opérant dans l'ensemble
d'indices $I$ de telle façon qu'on ait $g. U_i=U_{g. i}$ pour tout
$i\in I$ : on dira que ${\mathbf U}$ est un $G$-\textit{recouvrement}. Si $P$ est un
préfaisceau abélien sur $X$ dans lequel $G$ opère, alors le complexe
$C({\mathbf U},P)$ des cochaines de ${\mathbf U}$ à coéfficients dans $P$ peut
être considéré comme un complexe de $G$-groupes de façon évidente. Nous
poserons

\LayerTarget{tohoku:equation:5.5.3}
\[
\tag{5.5.3}
H^n({\mathbf U};G,P)={\mathbf R}^n\Gamma^G(C({\mathbf U},P))
\]
\LayerTarget{tohoku:construction:equivariant-cech-bicomplex}
où le deuxième membre peut donc aussi s'expliciter comme le groupe
$H^n(C(G,C({\mathbf U},P)))$, $C(G,C({\mathbf U},P))$ étant le bicomplexe formé
des cochaines de $G$ à valeurs dans le complexe $C({\mathbf U},P)$, bicomplexe
qu'on pourra noter $C({\mathbf U};G,P)$ :

\LayerTarget{tohoku:equation:5.5.3-bis}
\[
\tag{5.5.3 bis}
H^n({\mathbf U};G,P)=H^n(C({\mathbf U};G,P)).
\]

\LayerTarget{tohoku:construction:open-g-cover-spectral-functor}
Supposons d'abord le $G$-recouvrement ${\mathbf U}$ \textit{ouvert}. Nous allons appliquer
le lemme 5.5.1 en prenant les catégories ${\mathbf C}^{X(G)}$, ${\mathbf C}^G$ (catégorie des
$G$-modules), ${\mathbf C}$ (catégorie des groupes abéliens), et les foncteurs
$\Gamma_X$ et $\Gamma^G$, et posant enfin $\mathbf K(A)=C({\mathbf U},A)$.
$\Gamma_X$ transforme bien objets injectifs en objets
$\Gamma^G$-acycliques (corollaire à la prop. 5.1.3), et les conditions (i) et
(ii) du lemme sont satisfaites, (i) parce que ${\mathbf U}$ est ouvert, et (ii)
en vertu de ce qui a été dit dans 3.8 (deuxième alinéa). Il reste à expliciter
$R^q\mathbf K(A)$, le calcul est immédiat (et a d'ailleurs été fait dans 3.8)
on trouve $R^q\mathbf K(A)=C({\mathbf U},H^q(A))$, où $H^q(A)$ désigne le
préfaisceau $H^q(A)(V)=H^q(V,A)$. On a donc prouvé :

\LayerTarget{tohoku:theorem:5.5.3}
\textsc{Théorème 5.5.3.} \textit{Soit $X$ un espace muni d'un groupe
$G$ d'opérateurs, soit ${\mathbf U}=(U_i)$ un $G$-recouvrement ouvert de $X$.
Alors il existe sur la catégorie ${\mathbf C}^{X(G)}$ des $G$-faisceaux abéliens sur $X$,
un foncteur spectral cohomologique, aboutissant au foncteur gradué
$(H^n(X;G,A))$, et dont le terme initial est}

\LayerTarget{tohoku:equation:5.5.4}
\[
\tag{5.5.4}
E_2^{p,q}(A)=H^p({\mathbf U};G,H^q(A))
\]
\textit{(où $H^q(A)$ est le préfaisceau $H^q(A)(V)=H^q(V,A)$ sur $X$).}

\LayerTarget{tohoku:map:5.5.5}
On en conclut comme d'habitude des ``edge-homomorphisms''
\[
\tag{5.5.5}
H^n(U;G,A)\to H^n(X;G,A)
\]
\LayerTarget{tohoku:conclusion:theorem-5.5.3:p1}
et une suite exacte à 5 termes, montrant entre autres que \textit{l'homomorphisme}

\TohokuContinuation{209}
